\section{Quantum machine learning (Day 1, afternoon)}

\subsection{An introduction to quantum computing approaches for optimization problems}
\speaker{Marco Baioletti}{Universit\`a degli Studi di Perugia}
This pedagogical lecture introduced the two main quantum-computing paradigms
relevant for optimization: gate-based programmable devices and quantum
annealers. After reviewing qubits, superposition, entanglement, measurement,
elementary gates, and circuit execution on NISQ (noisy intermediate-scale
quantum) hardware, the speaker presented variational quantum algorithms, in
which a parametric circuit is optimized by a classical outer loop, with QAOA
(alternating problem and mixer Hamiltonians) as the prototypical algorithm for
combinatorial problems---illustrated with charged-particle track reconstruction
formulated as a QUBO over triplet-selection variables---and VQE for
Hamiltonian ground-state problems. The second part covered adiabatic quantum
computation and D-Wave-type annealers: the adiabatic theorem, annealing
schedules, quantum tunneling as a mechanism to escape local minima, the QUBO
formulation (with minimum vertex cover as a worked example), and practical
knobs such as number of reads, annealing time, and reverse annealing for hybrid
quantum--classical refinement. The take-away was that variational algorithms
attempt to work around NISQ limitations on gate-based machines, while annealers
are a pragmatic alternative for mid-size optimization problems.

\subsection{Quantum-Informed Neural Networks}
\speaker{Aritra Bal, Michael Binder}{KIT; with M.~Klute, B.~Maier, M.~Spannowsky}
Rather than running full quantum classifiers, this work extracts geometric
information from quantum representations of jets and injects it into classical
networks. Jets are encoded with the one-particle-one-qubit (1P1Q) scheme (up to
ten constituents, $(\pT,\eta,\phi)$ per qubit); the Quantum Fisher Information
Matrix (QFIM) of the resulting state is interpreted as a learned summary of
two-particle correlations: for top jets, large off-diagonal QFIM blocks connect
subjet constituents, while QCD jets show suppressed off-diagonals. Three
increasingly powerful applications were shown: (i) a training-free,
QFIM-only prototype classifier reaching VQC-level performance with faster,
more stable convergence when used as initialization; (ii) a non-trainable graph
classifier and a trainable GNN (``QuantumConv1D'') whose messages are augmented
with QFI sub-matrices, converging faster and outperforming the classical
benchmark; and (iii) QUIVER, which injects QFI tokens and attention biases into
the state-of-the-art Particle Transformer, giving statistically significant
improvements over five runs on top tagging and $H\to c\bar c$ benchmarks
(arXiv:2606.02785, accepted at the ICML~2026 AI4Physics workshop). Current
limitations are the $\sim$10-qubit simulability ceiling and the cost of
QFIM-based optimization; the outlook includes extending from top to boosted
$W/Z/H$ jets---directly relevant to polarization tagging---and hybrid
quantum-to-classical pipelines beyond qubits.

\subsection{Quantum Lund plane for VBS}
\speaker{Fabrizio Napolitano}{Universit\`a e INFN Perugia}
This talk (based on arXiv:2604.18613) proposed a quantum machine-learning
approach to boosted-boson and polarization tagging built on the Lund jet plane.
The LP2B (Lund Plane to Bloch) encoding maps each Lund declustering node,
through a differentiable and trainable stereographic projection, onto a qubit
of a Quantum Tree Tensor Network (QTTN) whose topology mirrors the C/A
declustering tree (kept shallow via a $\ln k_t > 1$ cut), with CRY
entanglement, variational layers, and readout; by construction the
representation is IRC-safe. Benchmarked against BDTs and MLPs on Lund inputs,
LundNet-2, and 1P1Q-style VQCs, on tasks including $W$ tagging, top tagging,
and $W$/graviton polarization discrimination, the QTTN reaches competitive
performance with a very small parameter count and no sign of overfitting even
with only $10^3$ training jets. A central physics message was that polarization
tagging is ``difficulty-limited'' ($W_T$ and $W_L$ are the same particle, so
methods saturate quickly), and that low-level constituent-based approaches
inherit parton-shower/hadronization modeling uncertainties of up to
$\sim$40\%, whereas IRC-safe high-level (Lund-plane) inputs are far more
resilient---a decisive consideration for analyses with large backgrounds in
the signal region. The model was validated on real quantum hardware (SpinQ
Triangulum, 3 qubits), with AUC compatible with simulation
($0.58\pm0.02$ vs $0.60\pm0.02$).

\subsection{Discussion}
The session discussion focused on where quantum approaches can be genuinely
complementary to classical ML---compact, physics-informed representations and
modeling robustness rather than raw performance---and on the scalability of
these methods towards realistic detector-level datasets.
